%%%%%%%%%%%%%%%%%%%%%%%%%%%%%%%%%%%%%%%%%
% Beamer Presentation
% LaTeX Template
% Version 1.0 (10/11/12)
%
% This template has been downloaded from:
% http://www.LaTeXTemplates.com
%
% License:
% CC BY-NC-SA 3.0 (http://creativecommons.org/licenses/by-nc-sa/3.0/)
%
%%%%%%%%%%%%%%%%%%%%%%%%%%%%%%%%%%%%%%%%%

%----------------------------------------------------------------------------------------
%	PACKAGES AND THEMES
%----------------------------------------------------------------------------------------

\documentclass[10pt]{beamer}

\mode<presentation> {

% The Beamer class comes with a number of default slide themes
% which change the colors and layouts of slides. Below this is a list
% of all the themes, uncomment each in turn to see what they look like.

%\usetheme{default}
%\usetheme{AnnArbor}
%\usetheme{Antibes}
%\usetheme{Bergen}
%\usetheme{Berkeley}
%\usetheme{Berlin}
%\usetheme{Boadilla}
%\usetheme{CambridgeUS}
%\usetheme{Copenhagen}
%\usetheme{Darmstadt}
%\usetheme{Dresden}
%\usetheme{Frankfurt}
%\usetheme{Goettingen}
%\usetheme{Hannover}
%\usetheme{Ilmenau}
%\usetheme{JuanLesPins}
%\usetheme{Luebeck}
%\usetheme{Madrid}
%\usetheme{Malmoe}
%\usetheme{Marburg}
%\usetheme{Montpellier}
%\usetheme{PaloAlto}
\usetheme{Pittsburgh}
%\usetheme{Rochester}
%\usetheme{Singapore}
%\usetheme{Szeged}
%\usetheme{Warsaw}

% As well as themes, the Beamer class has a number of color themes
% for any slide theme. Uncomment each of these in turn to see how it
% changes the colors of your current slide theme.

%\usecolortheme{albatross}
\usecolortheme{beaver}
%\usecolortheme{beetle}
%\usecolortheme{crane}
%\usecolortheme{dolphin}
%\usecolortheme{dove}
%\usecolortheme{fly}
%\usecolortheme{lily}
%\usecolortheme{orchid}
%\usecolortheme{rose}
%\usecolortheme{seagull}
%\usecolortheme{seahorse}
%\usecolortheme{whale}
%\usecolortheme{wolverine}

%\setbeamertemplate{footline} % To remove the footer line in all slides uncomment this line
%\setbeamertemplate{footline}[page number] % To replace the footer line in all slides with a simple slide count uncomment this line

%\setbeamertemplate{navigation symbols}{} % To remove the navigation symbols from the bottom of all slides uncomment this line
}

\usepackage{graphicx} % Allows including images
\usepackage{booktabs} % Allows the use of \toprule, \midrule and \bottomrule in tables
%\usepackage{ctex}
\usepackage{multicol}

\usepackage{fontspec}
\setsansfont{Calibri}[
    Path=Font/,
    Scale=0.9,
    Extension = .ttf,
    UprightFont=*-Regular,
    BoldFont=*-Bold,
    ItalicFont=*-Italic,
    BoldItalicFont=*-BoldItalic
    ]
\usepackage{unicode-math}

% \setmathfont{Libertinus Math} % 使用默认的数学字体或你偏好的数学字体 Latin Modern Math

\usepackage{ragged2e}
\justifying\let\raggedright\justifying %两端对齐
\usepackage{tikz}
\usetikzlibrary{positioning}
\usetikzlibrary{calc}
%----------------------------------------------------------------------------------------
%	TITLE PAGE
%----------------------------------------------------------------------------------------

\title[Asset Pricing]{\textbf{Asset Pricing}} % The short title appears at the bottom of every slide, the full title is only on the title page

\subtitle{Beta}

\author[Cheng Hang]{Cheng Hang} % Your name
\institute[DUFE] % Your institution as it will appear on the bottom of every slide, may be shorthand to save space
{School of Fintech \\
Dongbei University of Finance \& Economics \\ % Your institution for the title page
\medskip
\textbf{chenghangch@qq.com}% Your email address
}
\date{\today} % Date, can be changed to a custom date

	\AtBeginSection[]
{
	\begin{frame}{Content}
		\tableofcontents[sectionstyle=show/shaded,subsectionstyle=show/shaded/hide] %这几个参数我也不知道该如何准确地解释，反正它们最终的效果是突出显示当前章节，而其它章节都进行了淡化处理
		\addtocounter{framenumber}{-1}  %目录页不计算页码
	\end{frame}
}

%----------------------------------------------------------------------------------------
%	Begin
%----------------------------------------------------------------------------------------

\begin{document}

\begin{frame}
\titlepage % Print the title page as the first slide
\end{frame}

%\begin{frame}
%\frametitle{Overview} % Table of contents slide, comment this block out to remove it
%\tableofcontents % Throughout your presentation, if you choose to use \section{} and \subsection{} commands, these will automatically be printed on this slide as an overview of your presentation
%\end{frame}

%---------------------------------------------------------------------------
%	PRESENTATION SLIDES
%---------------------------------------------------------------------------
\section{Fama and MacBeth (1973)}

% --- Introduction ---
\begin{frame}{The Core Question}
  \begin{itemize}
    \item Empirical asset pricing tests aim to determine if factors have significant risk premia.
    \item A crucial first step: Correctly calculating factor loadings ($\beta$).
  \end{itemize}
\end{frame}

% --- Fama-MacBeth Regression ---
\begin{frame}{Background: Fama-MacBeth (1973) Regression}
  A standard two-step procedure to test for risk premia:
  \begin{enumerate}
    \item \textbf{First Pass (Time-Series Regression):}
      \begin{itemize}
        \item For each asset $i$, regress its returns $R_{i,t}$ on factor returns $F_{k,t}$ over time:
              \[ R_{i,t} = \alpha_i + \sum_{k=1}^{K} \beta_{i,k} F_{k,t} + \epsilon_{i,t} \]
        \item This yields estimates of factor loadings $\hat{\beta}_{i,k}$ for each asset $i$ on each factor $k$.
      \end{itemize}
    \item \textbf{Second Pass (Cross-Sectional Regression):}
      \begin{itemize}
        \item For each time period $t$, regress asset returns $R_{i,t}$ on their estimated factor loadings $\hat{\beta}_{i,k}$ from the first pass:
              \[ R_{i,t} = \lambda_{0,t} + \sum_{k=1}^{K} \lambda_{k,t} \hat{\beta}_{i,k} + \eta_{i,t} \]
        \item This yields estimates of risk premia $\hat{\lambda}_{k,t}$ for each factor $k$ at time $t$.
        \item The final risk premium for factor $k$ is the time-series average of $\hat{\lambda}_{k,t}$, and its significance is tested (e.g., t-test).
      \end{itemize}
  \end{enumerate}
\end{frame}

% --- Theoretical Background ---
\begin{frame}
  \frametitle{Theoretical Underpinnings: The Two-Parameter Model}
  The paper tests three key implications of the traditional CAPM:
  \begin{enumerate}
    \item \textbf{Linearity:} The relationship between an asset's expected return and its systematic risk (beta) is linear.
    \item \textbf{Systematic Risk is Key:} No other measure of risk, particularly non-systematic (idiosyncratic) risk, systematically affects expected returns.
    \item \textbf{Positive Risk-Return Tradeoff:} There is a positive expected premium for bearing systematic risk.
  \end{enumerate}

  \bigskip
  The ex-ante CAPM equation:
  \[ E(R_i) = R_f + \beta_i [E(R_m) - R_f] \]
  where:
  \begin{itemize}
    \item $E(R_i)$: Expected return on asset $i$.
    \item $R_f$: Risk-free rate of return.
    \item $\beta_i = \text{Cov}(R_i, R_m) / \text{Var}(R_m)$: Systematic risk of asset $i$.
    \item $E(R_m)$: Expected return on the market portfolio.
  \end{itemize}
  Fama and MacBeth aimed to test a stochastic (ex-post) version of this model.
\end{frame}

% --- Data and Period ---
\begin{frame}
  \frametitle{Data and Sample Period}
  \begin{itemize}
    \item \textbf{Assets:} All common stocks listed on the New York Stock Exchange (NYSE).
    \item \textbf{Period:} Monthly data from January 1926 to June 1968.
      \begin{itemize}
        \item Overall period: 1935-1968 (for the main tests).
        \item Subperiods were also examined: 1935-1945, 1946-1955, 1956-1968.
      \end{itemize}
    \item \textbf{Market Proxy ($R_m$):} An equally-weighted portfolio of all NYSE stocks.
    \item \textbf{Risk-Free Rate ($R_f$):} One-month Treasury bill rate (or commercial paper rates for earlier periods where T-bills were not available).
  \end{itemize}
\end{frame}

% --- Methodology Overview ---
\begin{frame}
  \frametitle{Methodology: The Two-Pass Regression Approach}
  The core of the Fama-MacBeth methodology involves a two-stage (or "two-pass") procedure repeated over time:

  \begin{block}{Overall Procedure}
  \begin{enumerate}
    \item \textbf{Portfolio Formation and Beta Estimation (First Pass - Time Series):}
        \begin{itemize}
            \item Form portfolios of stocks to mitigate the errors-in-variables (EIV) problem associated with individual stock beta estimates.
            \item Estimate betas for these portfolios using time-series regressions.
        \end{itemize}
    \item \textbf{Cross-Sectional Regressions (Second Pass):}
        \begin{itemize}
            \item For each month in the sample period, run a cross-sectional regression of portfolio returns on their estimated betas (and other test variables).
        \end{itemize}
    \item \textbf{Hypothesis Testing:}
        \begin{itemize}
            \item Analyze the time series of coefficients obtained from the monthly cross-sectional regressions.
        \end{itemize}
  \end{enumerate}
  \end{block}
\end{frame}

% --- Step 1: Portfolio Formation and Beta Estimation ---
\begin{frame}[allowframebreaks]
  \frametitle{Step 1 (First Pass): Portfolio Formation \& Beta Estimation}
  \textbf{Objective:} Obtain reliable estimates of systematic risk ($\beta$) and ensure a wide dispersion of betas across test assets.

  \begin{enumerate}
    \item \textbf{Initial Beta Ranking Period (e.g., 1926-1929 for first portfolio set):}
      \begin{itemize}
        \item For each individual stock, estimate its beta ($\hat{\beta}_i$) using monthly returns over a preceding period (e.g., 4-7 years, typically 5 years in their main setup) by regressing individual stock returns on the market proxy returns:
              \[ R_{i,t} = a_i + b_i R_{m,t} + e_{i,t} \]
              (Note: They used raw returns, not excess returns, for this initial beta estimation.)
      \end{itemize}

    \item \textbf{Portfolio Formation (e.g., at end of 1929):}
      \begin{itemize}
        \item Rank all stocks based on their estimated $\hat{\beta}_i$ from the ranking period.
        \item Group these stocks into 20 portfolios (equally weighted returns for each portfolio). This helps diversify away idiosyncratic risk and provides more stable beta estimates for the portfolios.
      \end{itemize}

    \item \textbf{Portfolio Beta Estimation Period (e.g., 1930-1934 for betas used in 1935 CSR):}
      \begin{itemize}
        \item For each of the 20 portfolios, calculate its monthly returns.
        \item Estimate the beta for each portfolio ($\hat{\beta}_{p,t-1}$) using monthly returns from a subsequent period (e.g., the next 5 years) by regressing portfolio returns on market proxy returns:
              \[ R_{p,\tau} = a_p + \beta_p R_{m,\tau} + e_{p,\tau} \]
              These $\hat{\beta}_{p,t-1}$ are the betas used as independent variables in the cross-sectional regressions for month $t$.
      \end{itemize}

    \item \textbf{Rolling Procedure:}
      \begin{itemize}
        \item The entire process (ranking, portfolio formation, portfolio beta estimation) is repeated annually. For example, to get betas for 1936, portfolios are reformed based on 1927-1930 individual stock betas, and portfolio betas are estimated using 1931-1935 data.
        \item This ensures that the betas used in the cross-sectional regressions are estimated from data prior to the month of the cross-sectional regression.
      \end{itemize}
  \end{enumerate}
\end{frame}

% --- Step 2: Cross-Sectional Regressions ---
\begin{frame}[allowframebreaks]
  \frametitle{Step 2 (Second Pass): Monthly Cross-Sectional Regressions}
  For \textbf{each month $t$} in the test period (e.g., January 1935 to June 1968), run the following cross-sectional regression:

  \[ R_{p,t} = \hat{\gamma}_{0,t} + \hat{\gamma}_{1,t} \hat{\beta}_{p,t-1} + \hat{\gamma}_{2,t} \hat{\beta}_{p,t-1}^2 + \hat{\gamma}_{3,t} \bar{s}_{p,t-1}(\hat{e}_i) + \eta_{p,t} \]
  where:
  \begin{itemize}
    \item $R_{p,t}$: The return on portfolio $p$ in month $t$. (They used both raw returns and excess returns $R_{p,t} - R_{f,t}$ in different specifications).
    \item $\hat{\beta}_{p,t-1}$: The estimated beta for portfolio $p$, calculated using data from the period \textit{ending before month $t$}.
    \item $\hat{\beta}_{p,t-1}^2$: The square of the portfolio's beta. This term tests for non-linearity in the risk-return relationship.
    \item $\bar{s}_{p,t-1}(\hat{e}_i)$: The average residual standard deviation of the stocks within portfolio $p$. This is a measure of non-systematic (idiosyncratic) risk, estimated from the first-pass individual stock regressions used for portfolio formation. This term tests if non-market risk is priced.
    \item $\hat{\gamma}_{0,t}, \hat{\gamma}_{1,t}, \hat{\gamma}_{2,t}, \hat{\gamma}_{3,t}$: The regression coefficients estimated for month $t$.
    \item $\eta_{p,t}$: The regression residual for portfolio $p$ in month $t$.
  \end{itemize}
  This regression is run for each of the $N_t$ months in the sample, yielding $N_t$ estimates for each $\gamma_j$.
\end{frame}

% --- Hypothesis Testing ---
\begin{frame}[allowframebreaks]
  \frametitle{Hypothesis Testing: Analyzing the Time Series of Coefficients}
  The CAPM implies specific hypotheses about the expected values of the $\gamma_j$ coefficients:

  \begin{enumerate}
    \item \textbf{Average Coefficients:} Calculate the time-series average of each estimated coefficient:
          \[ \bar{\hat{\gamma}}_j = \frac{1}{N_t} \sum_{t=1}^{N_t} \hat{\gamma}_{j,t} \quad \text{for } j = 0, 1, 2, 3 \]

    \item \textbf{Test Statistics:} Test if the average coefficients are significantly different from their hypothesized values using a t-statistic:
          \[ t(\bar{\hat{\gamma}}_j) = \frac{\bar{\hat{\gamma}}_j - E(\gamma_j)_{\text{hypothesized}}}{s(\hat{\gamma}_j) / \sqrt{N_t}} \]
          where $s(\hat{\gamma}_j)$ is the standard deviation of the time series of $\hat{\gamma}_{j,t}$.

    \item \textbf{Specific Hypotheses (based on $R_{p,t}$ as LHS):}
      \begin{itemize}
        \item $E(\gamma_{0,t})$: Expected return on a zero-beta portfolio ($E(R_z)$). If the strict CAPM holds, $E(\gamma_{0,t})$ should be $R_{f,t}$ (or its average).
        \item $E(\gamma_{1,t})$: Expected market risk premium ($E(R_m) - E(R_z)$ or $E(R_m) - R_f$). Should be positive.
        \item $E(\gamma_{2,t})$: Tests for non-linearity. Should be zero.
        \item $E(\gamma_{3,t})$: Tests if non-systematic risk is priced. Should be zero.
      \end{itemize}
      (If $R_{p,t} - R_{f,t}$ is LHS, then $E(\gamma_{0,t})$ should be zero.)
  \end{enumerate}
  The significance of these t-statistics provides evidence for or against the CAPM's predictions.
\end{frame}

% --- Key Findings (Brief Summary) ---
\begin{frame}[allowframebreaks]
  \frametitle{Summary of Fama-MacBeth (1973) Key Findings}
  While the paper is primarily about methodology, its main empirical findings were:

  \begin{itemize}
    \item \textbf{Support for Linearity and Systematic Risk:}
      \begin{itemize}
        \item The average estimate $\bar{\hat{\gamma}}_2$ (for $\beta^2$) was generally not significantly different from zero, supporting linearity.
        \item The average estimate $\bar{\hat{\gamma}}_3$ (for residual risk $\bar{s}(\hat{e}_i)$) was generally not significantly different from zero, supporting that only systematic risk is priced.
      \end{itemize}
    \item \textbf{Positive Risk-Return Tradeoff:}
      \begin{itemize}
        \item The average estimate $\bar{\hat{\gamma}}_1$ (the coefficient on $\beta$) was consistently positive and statistically significant, indicating a positive price for systematic risk.
      \end{itemize}
    \item \textbf{Deviations from Strict CAPM (Sharpe-Lintner version):}
      \begin{itemize}
        \item $\bar{\hat{\gamma}}_0$ was often significantly higher than the average risk-free rate.
        \item $\bar{\hat{\gamma}}_1$ was often significantly lower than the average market excess return ($R_m - R_f$).
        \item These findings were more consistent with Black's (1972) zero-beta CAPM, where $E(\gamma_{0,t}) = E(R_{z,t})$ and $E(\gamma_{1,t}) = E(R_{m,t}) - E(R_{z,t})$, and $E(R_{z,t})$ could be greater than $R_{f,t}$ due to restrictions on borrowing.
      \end{itemize}
    \item \textbf{Stability:} The qualitative results were generally stable across different subperiods.
  \end{itemize}
\end{frame}

% --- Contribution and Legacy ---
\begin{frame}
  \frametitle{Contribution and Legacy of Fama-MacBeth (1973)}
  \begin{itemize}
    \item \textbf{Standard Methodology:} The Fama-MacBeth two-pass regression procedure became a standard tool for testing asset pricing models and examining the cross-section of stock returns.
    \item \textbf{Addressing EIV:} The use of portfolios was a key innovation to mitigate the errors-in-variables problem in beta estimation.
    \item \textbf{Time-Varying Risk Premia:} By estimating coefficients for each month, the methodology implicitly allows for time-varying risk premia. The test is on the average premium.
    \item \textbf{Foundation for Future Research:} It laid the groundwork for countless subsequent studies on asset pricing anomalies and factor models.
    \item \textbf{Limitations Acknowledged:}
        \begin{itemize}
            \item Market proxy issues (Roll's Critique).
            \item Assumes betas are relatively stable within estimation periods.
        \end{itemize}
  \end{itemize}
\end{frame}

\section{Standard Estimation beta in recent Literature}

\begin{frame}{Method 1: Standard CAPM Regression}
  The most common approach is the one-factor market model regression:
  \begin{equation*}
    r_{i,t} = \alpha_i + \beta_i \text{MKT}_t + \epsilon_{i,t} \quad (8.3)
  \end{equation*}
  \begin{itemize}
    \item $r_{i,t}$: Excess return of stock $i$ ($R_{i,t} - R_{f,t}$)
    \item $\text{MKT}_t$: Excess return of the market portfolio ($R_{m,t} - R_{f,t}$)
    \item $\epsilon_{i,t}$: Regression residual
    \item The estimated slope coefficient is taken as the stock's market beta ($\hat{\beta}_i$).
  \end{itemize}

  \textbf{Key Variations:}
  \begin{itemize}
    \item Periodicity of data (daily, monthly)
    \item Length of the estimation period
  \end{itemize}
\end{frame}

\begin{frame}{Standard CAPM Regression: Data Variations}
  \textbf{Using Daily Excess Return Data:}
  \begin{itemize}
    \item Notation: $\beta^{kM}$, where $k$ is the number of months.
    \item Examples from the text:
    \begin{itemize}
        \item $\beta^{1M}$: 1 month of daily data (min 15 days)
        \item $\beta^{3M}$: 3 months of daily data (min 50 days)
        \item $\beta^{6M}$: 6 months of daily data (min 100 days)
        \item $\beta^{12M}$: 12 months of daily data (min 200 days)
        \item $\beta^{24M}$: 24 months of daily data (min 450 days)
    \end{itemize}
  \end{itemize}

  \textbf{Using Monthly Excess Return Data:}
  \begin{itemize}
    \item Notation: $\beta^{kY}$, where $k$ is the number of years.
    \item Examples from the text:
    \begin{itemize}
        \item $\beta^{1Y}$: 1 year of monthly data (min 10 obs.)
        \item $\beta^{2Y}$: 2 years of monthly data (min 20 obs.)
        \item $\beta^{3Y}$: 3 years of monthly data (min 24 obs.)
        \item $\beta^{5Y}$: 5 years of monthly data (min 24 obs.)
    \end{itemize}
  \end{itemize}
\end{frame}

\begin{frame}{Method 2: Scholes-Williams (1977) Beta}
  \textbf{Purpose:} To account for nonsynchronous trading, which can affect beta estimates from the standard CAPM model, especially with daily data.

  \textbf{Procedure:} Run three regressions:
  \begin{align*}
    r_{i,t} &= a_i + b_i^{-} \text{MKT}_{t-1} + e_{i,t}^{-} \\
    r_{i,t} &= a_i + b_i \text{MKT}_{t} + e_{i,t}  \\
    r_{i,t} &= a_i + b_i^{+} \text{MKT}_{t+1} + e_{i,t}^{+}
  \end{align*}

  The Scholes-Williams beta is defined as:
  \begin{equation*}
    \beta_i^{\text{SW}} = \frac{\hat{b}_i^{-} + \hat{b}_i + \hat{b}_i^{+}}{1 + 2\rho} 
  \end{equation*}
  where $\rho$ is the first-order serial correlation of the market portfolio's excess return.

  \textbf{Implementation (from text):}
  \begin{itemize}
    \item Calculated using daily data.
    \item Example: 1-year period (months $t-11$ to $t$), min 200 observations.
  \end{itemize}
\end{frame}

\begin{frame}{Method 3: Dimson (1979) Beta}
  \textbf{Purpose:} To address bias in beta estimates for infrequently traded stocks when using the standard CAPM model.\\

  \textbf{Procedure:} Estimate coefficients from the following regression:
  \begin{equation*}
    r_{i,t} = a_i + \sum_{k=-L}^{L} b_{k,i} \text{MKT}_{t+k} + e_{i,t} 
  \end{equation*}
  The text uses $L=5$:
  \begin{equation*}
    r_{i,t} = a_i + \sum_{k=-5}^{5} b_{k,i} \text{MKT}_{t+k} + e_{i,t} 
  \end{equation*}

  The Dimson beta is the sum of these slope coefficients:
  \begin{equation*}
    \beta_i^{\text{D}} = \sum_{k=-5}^{5} \hat{b}_{k,i} 
  \end{equation*}

  \textbf{Implementation (from text):}
  \begin{itemize}
    \item Applicable to all stocks (does not introduce bias for frequently traded ones).
    \item Example: 1-year daily data (months $t-11$ to $t$), min 200 observations.
  \end{itemize}
\end{frame}

\begin{frame}{Data Sources and Sample (as per the text)}
  The empirical examination in the chapter uses:
  \begin{itemize}
    \item \textbf{Stock Return Data:}
    \begin{itemize}
        \item Daily: Center for Research in Security Prices (CRSP) daily stock file (dsf).
        \item Monthly: CRSP monthly stock file (msf).
    \end{itemize}
    \item \textbf{Market Factor (MKT) and Risk-Free Rate ($R_f$):}
    \begin{itemize}
        \item Daily and Monthly: Ken French's data library.
        \item Used to calculate excess returns.
    \end{itemize}
    \item \textbf{Sample Period for Empirical Examination:}
    \begin{itemize}
        \item June 1963 through November 2012.
    \end{itemize}
  \end{itemize}
  \textit{Note: These are details in Chapter in ``Book Bali, Turan, Robert Engle, and Scott Murray, 2017, Empirical Asset Pricing: The Cross Section of Stock Returns: An Overview".
.}
\end{frame}

\begin{frame}{Summary of Beta Estimation Methods}
  This chapter excerpt discusses three main approaches to estimating market beta:
  \begin{enumerate}
    \item \textbf{Standard CAPM Regression:}
    \begin{itemize}
        \item $r_{i,t} = \alpha_i + \beta_i \text{MKT}_t + \epsilon_{i,t}$
        \item Variations based on data frequency (daily/monthly) and estimation window length.
    \end{itemize}
    \item \textbf{Scholes-Williams (1977) Beta:}
    \begin{itemize}
        \item Adjusts for nonsynchronous trading using lead, contemporaneous, and lagged market returns.
        \item $\beta_i^{\text{SW}} = (\hat{b}_i^{-} + \hat{b}_i + \hat{b}_i^{+}) / (1 + 2\rho)$
    \end{itemize}
    \item \textbf{Dimson (1979) Beta:}
    \begin{itemize}
        \item Addresses infrequent trading by summing coefficients from a regression with multiple lead and lagged market returns.
        \item $\beta_i^{\text{D}} = \sum \hat{b}_{k,i}$
    \end{itemize}
  \end{enumerate}
  These methods are widely used in empirical asset pricing research.
\end{frame}

\section{EIV}

% --- The EIV Problem ---

\begin{frame}{The Errors-in-Variables (EIV) Problem}
  \begin{itemize}
    \item Using $\hat{\beta}_{i,k}$ instead of true $\beta_{i,k}$ leads to an Errors-in-Variables (EIV) bias.
    \item Risk premium estimates ($\hat{\gamma}_k$) become biased and inconsistent.
    \item Direction of bias is unknown with multiple factors.
    \item \textbf{Traditional Solution:} Use portfolios as test assets instead of individual stocks.
    \begin{itemize}
        \item Portfolio betas are estimated more precisely, reducing EIV bias.
    \end{itemize}
  \end{itemize}
\end{frame}

\subsection{Jegadeesh et al. (2019)}

\begin{frame}[allowframebreaks]{Motivation: Shortcomings of Using Portfolios}
  \textbf{Why move away from portfolios? Several critical issues:}
  \begin{itemize}
    \item \textbf{Reduced Test Power:} Dimensionality is reduced. Average returns vary less across portfolios than individual stocks.
    \item \textbf{Masking Phenomena:} Diversification into portfolios can hide cross-sectional patterns in individual stocks unrelated to the grouping procedure (e.g., fundamental indexation arguments).
    \item \textbf{Misleading Inference on Efficiency:}
        \begin{itemize}
            \item If the market index is on the mean-variance frontier of individual assets, a linear relation between E(R) and $\beta$ holds.
            \item Grouping by beta might lead to the mistaken inference that the index is efficient with respect to beta-sorted portfolios, even if not for individual assets.
        \end{itemize}
    \item \textbf{Factor Structure from Sorting (Lewellen et al., 2010):}
        \begin{itemize}
            \item Sorting on characteristics (size, B/M) known to predict returns imparts a strong factor structure in test portfolio returns.
            \item Even factors weakly correlated with sorting characteristics might appear to explain average returns.
        \end{itemize}
    \item \textbf{Dependence on Test Portfolios:} Significance and magnitude of risk premiums can depend critically on the choice of test portfolios.
  \end{itemize}
\end{frame}

\begin{frame}{This Paper's Contribution}
  \begin{itemize}
    \item \textbf{Objective:} Develop a procedure to estimate risk premiums and test their significance using \textbf{individual stocks} while avoiding EIV bias.
    \item \textbf{Proposed Method:} An Instrumental Variables (IV) approach.
    \begin{itemize}
        \item Allows use of individual stocks as test assets.
        \item Delivers consistent estimates of ex-post risk premiums (N-consistency).
        \item Yields well-specified tests in small samples.
    \end{itemize}
    \item \textbf{Key Punchline (Preview):} Market risk premium (CAPM, LCAPM), premiums on FF3, FF5, and Hou et al. (2015) q-factors are all \textbf{insignificant} after controlling for asset characteristics when using this IV method with individual stocks.
  \end{itemize}
\end{frame}

\begin{frame}{The IV Estimator}
  Recall the second-pass regression for month $t$:
  \[ r_{i,t} = \gamma_{0,t} + \sum_{k=1}^{K} \hat{\beta}_{i,k} \gamma_{k,t} + \xi_{i,t} \]
  In matrix form: $r_t = \hat{\gamma} \hat{B} + \xi_t$.

  The proposed IV estimator for factor risk premiums $\hat{\gamma}_{IV,t}$ is:
  \[ \hat{\gamma}'_{IV,t} = (\hat{B}_{IV} \hat{B}'_{EV})^{-1} (\hat{B}_{IV} r'_t) \]
  Where:
  \begin{itemize}
    \item $\hat{B}_{EV}$: Matrix of explanatory variable betas (e.g., estimated using even months' daily data from a prior period).
    \item $\hat{B}_{IV}$: Matrix of instrumental variable betas (e.g., estimated using odd months' daily data from the same prior period).
    \item $r_t$: Vector of realized excess returns in month $t$.
  \end{itemize}

  \textbf{Key Idea:} Measurement errors in $\hat{B}_{EV}$ and $\hat{B}_{IV}$ are uncorrelated because they are estimated from disjoint data samples.
\end{frame}

% \begin{frame}{Properties of the IV Estimator}
%   \begin{itemize}
%     \item \textbf{N-Consistency (Proposition 1):} The IV estimator converges to the ex-post risk premium as the number of stocks (N) increases, even for finite T.
%     \item \textbf{Intuition (2SLS):}
%     \begin{enumerate}
%         \item Regress explanatory betas ($\hat{B}_{EV}$) on instrumental betas ($\hat{B}_{IV}$).
%         \item Use fitted values from (1) as explanatory variables in the second-stage regression for $r_t$.
%     \end{enumerate}
%     \item \textbf{Small Sample Properties (Simulations):}
%     \begin{itemize}
%         \item IV estimates are not different from true parameters even with short time-series for beta estimation (e.g., 1 year of daily data).
%         \item OLS estimates with individual stocks are significantly biased towards zero.
%         \item Conventional t-tests based on IV estimator are well-specified (correct size) and reasonably powerful.
%     \end{itemize}
%     \item \textbf{Time-Varying Betas:} IV estimator remains consistent if betas follow covariance stationary processes.
%   \end{itemize}
% \end{frame}

% \begin{frame}{Data and General Setup}
%   \begin{itemize}
%     \item \textbf{Data:} CRSP (returns, volume, market cap), Compustat (financials).
%     \item \textbf{Sample Period:} Jan 1956 - Dec 2012 (main tests for CAPM/FF3). Varies for other models based on factor availability.
%     \item \textbf{Stock Inclusion:} Common stocks (CRSP share codes 10/11), price > \$1, market cap > \$1M. Sufficient daily return data for beta estimation.
%     \item \textbf{Beta Estimation:}
%         \begin{itemize}
%             \item Daily returns over rolling 36-month windows.
%             \item Dimson (1979) betas (with 1-day lead/lag) to account for non-synchronous trading.
%             \item For IV: Odd-month daily data for one set of betas, even-month daily data for the other.
%         \end{itemize}
%     \item \textbf{Characteristics:} SIZE (log market cap), BM (book-to-market), OP (operating profitability), INV (investment).
%   \end{itemize}
% \end{frame}

\begin{frame}{Summary of Main Empirical Findings}
  \begin{itemize}
    \item Across multiple prominent asset pricing models (CAPM, FF3F, FF5F, q-Factor):
    \item When using \textbf{individual stocks} as test assets and the proposed \textbf{IV method}:
    \item \textbf{Factor risk premiums are generally NOT statistically significant after controlling for corresponding firm characteristics.}
    \item \textbf{Firm characteristics themselves (SIZE, BM, OP, INV, Amihud Illiquidity) are strongly and significantly priced.}
    \item This contrasts with findings from original papers that typically use portfolios and find significant risk premiums.
    \item Simulation evidence suggests this is not due to a lack of test power of the IV method.
  \end{itemize}
\end{frame}

\begin{frame}{The ``Risk-Proxy" Hypothesis}
  \begin{itemize}
    \item \textbf{Concern:} Could characteristics be priced because they are better proxies for \textit{true (future) betas} than betas estimated from past data, especially if true betas are time-varying and characteristics anticipate these changes?
    \item If so, the significant coefficients on characteristics might actually reflect compensation for systematic risk.
  \end{itemize}
\end{frame}

\section{Downside Beta}

\subsection{Ang, Chen, and Xing (2006)}

%----------------------------------------------------------------------
\begin{frame}{Introduction \& Motivation}
    \begin{itemize}
        \item Economists have long recognized that \textbf{investors care differently about downside losses versus upside gains}.
        \item Standard CAPM beta is symmetric and may not fully capture risk.
        \item This paper investigates whether \textbf{stocks with high sensitivities to downside market movements (high ``downside risk") earn a return premium}.
        \item Key Question: Is downside risk priced in the cross-section of stock returns?
    \end{itemize}
\end{frame}

%----------------------------------------------------------------------
\begin{frame}{What is Downside Risk?}
    \begin{columns}[T] % Align columns at the top
        \begin{column}{0.5\textwidth}
            \textbf{Conceptual Understanding}
            \begin{itemize}
                \item Assets that perform poorly when the overall market declines are particularly undesirable.
                \item Investors require compensation for holding such assets.
            \end{itemize}
            \vspace{1em}
            \textbf{Measuring Downside Beta ($\beta^{-}$)}
            \begin{itemize}
                \item Defined analogously to CAPM beta, but conditional on market declines:
                \end{itemize}
                $$
 \beta_i^{-} = \frac{Cov(R_i, R_m | R_m < \mu_m)}{Var(R_m | R_m < \mu_m)} 
$$
                 where $R_m < \mu_m$ indicates periods where excess market return is below its mean.
        \end{column}
        \begin{column}{0.5\textwidth}
            \textbf{Other Related Measures}
            \begin{itemize}
                \item \textbf{Upside Beta ($\beta^{+}$):} Covariance when $R_m > \mu_m$.
                \item \textbf{Relative Downside Beta:} $\beta^{-} - \beta_{CAPM}$.
                \item This paper focuses on the idea that high $\beta^{-}$ should command a premium.
            \end{itemize}
        \end{column}
    \end{columns}
\end{frame}

\begin{frame}{Theoretical Motivation: Disappointment Aversion}
    \begin{itemize}
        \item The paper uses a simple model with a representative agent exhibiting \textbf{Disappointment Aversion (DA)} preferences (Gul, 1991).
        \item DA utility places greater weight on outcomes below a certain endogenous reference point (the certainty equivalent of wealth) than on outcomes above it.
        \item This framework naturally leads to a higher required return for assets that covary more strongly with the market during market downturns.
        \item \textbf{The model predicts that assets with higher downside risk ($\beta^{-}$) will have higher expected returns, not fully captured by standard CAPM beta.}
        \item It also implies a potential discount for assets with high upside potential (high $\beta^{+}$).
    \end{itemize}
\end{frame}

\begin{frame}{Contemporaneous Premium: Portfolio Sorts}
    \begin{itemize}
        \item Strategy: Sort stocks into quintiles based on their \textit{realized} betas (e.g., $\beta^{-}$, relative $\beta^{-}$) computed over a 12-month period, and examine contemporaneous returns of these portfolios.
        \item Data: NYSE stocks, July 1963 - December 2001. Daily data for beta calculations.
        \item \textbf{Key Findings (from Table 1 in the paper):}
        \begin{itemize}
            \item Stocks with high realized $\beta^{-}$ earn significantly higher average returns.
                \begin{itemize}
                    \item \textbf{High-Low $\beta^{-}$ portfolio spread: 11.8\% per annum.}
                \end{itemize}
            \item Stocks with high realized \textit{relative} $\beta^{-}$ (i.e., $\beta^{-} - \beta_{CAPM}$) also earn high returns.
                 \begin{itemize}
                    \item \textbf{High-Low relative $\beta^{-}$ spread: 6.6\% per annum.}
                \end{itemize}
            \item Stocks with high realized \textit{relative} $\beta^{+}$ (i.e., $\beta^{+} - \beta_{CAPM}$) earn lower returns, suggesting a discount for upside potential.
                \begin{itemize}
                    \item \textbf{High-Low relative $\beta^{+}$ spread: -6.1\% per annum.}
                \end{itemize}
        \end{itemize}
    \end{itemize}  
\end{frame}

\begin{frame}{Contemporaneous Premium: Fama-MacBeth Regressions}
    \begin{itemize}
        \item Fama-MacBeth (1973) cross-sectional regressions of monthly excess returns on firm characteristics and realized betas.
        \item Controls include: size, book-to-market, momentum, volatility, coskewness, cokurtosis, and liquidity beta.
        \item \textbf{Key Findings (from Table 2 in the paper):}
        \begin{itemize}
            \item The coefficient on \textbf{$\beta^{-}$ is consistently positive and statistically significant.}
            \item Estimated premium for downside risk is approximately \textbf{6\% per annum} after various controls.
            \item The coefficient on $\beta^{+}$ is often negative but less robust than the $\beta^{-}$ premium.
            \item The downside risk premium is distinct from and incremental to other known cross-sectional effects like size, value, momentum, and coskewness.
        \end{itemize}
    \end{itemize}
\end{frame}
%----------------------------------------------------------------------
% \begin{frame}{Downside Beta vs. Coskewness}
%     \begin{itemize}
%         \item Coskewness (Harvey and Siddique, 2000) also captures asymmetry, where investors prefer positive coskewness (dislike negative coskewness).
%         \item Are downside beta and coskewness capturing the same thing?
%         \item \textbf{Double Sorts (Table 3 in the paper):}
%         \begin{itemize}
%             \item First sort on coskewness, then on $\beta^{-}$ (and vice-versa).
%             \item The \textbf{premium for high $\beta^{-}$ persists even after controlling for coskewness levels.} (High-Low $\beta^{-}$ spread of 7.6\% p.a. within coskewness quintiles).
%             \item The \textbf{premium for low coskewness persists after controlling for $\beta^{-}$ levels.} (High-Low coskewness spread of -6.2\% p.a. within $\beta^{-}$ quintiles).
%         \end{itemize}
%         \item Conclusion: **Downside beta and coskewness are related but distinct risk factors.** Downside beta explicitly conditions on market declines.
%     \end{itemize}
% \end{frame}
% %----------------------------------------------------------------------
% \begin{frame}{Robustness of Contemporaneous Premium}
%     The paper performs extensive robustness checks for the contemporaneous downside risk premium. The findings are generally robust to:
%     \begin{itemize}
%         \item \textbf{Different Cutoff Points for Beta Calculation:} Using risk-free rate or zero return as market benchmark instead of mean market return.
%         \item \textbf{Data Frequency:} Using weekly data over 24 months for beta calculation.
%         \item \textbf{Portfolio Weighting:} Value-weighted portfolios (premium slightly smaller but still significant).
%         \item \textbf{Stock Universe:} Including all NYSE/AMEX/NASDAQ stocks (premium tends to be larger).
%         \item \textbf{Overlapping vs. Non-overlapping Periods:} Results hold for non-overlapping annual periods.
%         \item \textbf{Downside Risk in Fama-French Portfolios (Table 6):} GMM tests show that downside market risk significantly helps price the 25 Fama-French size and book-to-market sorted portfolios.
%     \end{itemize}
% \end{frame}

% \begin{frame}{Predicting Future Downside Risk Exposure}
%     For a downside risk premium to be exploitable, past downside risk exposure should predict future exposure.
%     \begin{itemize}
%         \item \textbf{Determinants of Future Relative $\beta^{-}$ (Table 7):}
%         \begin{itemize}
%             \item Past relative $\beta^{-}$ weakly predicts future relative $\beta^{-}$ (12-month autocorrelation is only 0.082).
%             \item Other firm characteristics like past standard deviation (+), log-size (-), and past ROE (+) show some predictive power for future relative $\beta^{-}$.
%         \end{itemize}
%          \item This suggests that simply using past $\beta^{-}$ might be a noisy predictor of future $\beta^{-}$.
%     \end{itemize}
% \end{frame}
% %----------------------------------------------------------------------
% \begin{frame}{Past Downside Risk and Future Returns}
%     \begin{itemize}
%         \item \textbf{Portfolio Sorts on Past $\beta^{-}$ (Table 8, Panel A):}
%         \begin{itemize}
%             \item Sorting stocks based on $\beta^{-}$ from the previous 12 months.
%             \item Result: A \textbf{weak and inconsistent predictive relationship between past $\beta^{-}$ and next month's returns.} The highest past $\beta^{-}$ quintile does not consistently earn the highest future returns.
%         \end{itemize}
%         \item \textbf{Contrast with Past Coskewness (Table 8, Panel B):}
%         \begin{itemize}
%             \item Past coskewness strongly predicts future returns (low past coskewness $\rightarrow$ high future returns), consistent with Harvey and Siddique (2000).
%         \end{itemize}
%         \item Why is the predictive power of past $\beta^{-}$ weak? \textbf{The role of volatility.}
%     \end{itemize}
% \end{frame}
% %----------------------------------------------------------------------
% \begin{frame}{The Confounding Role of Volatility}
%     Two issues with high volatility stocks:
%     \begin{enumerate}
%         \item \textbf{Measurement Error:} High volatility makes past $\beta^{-}$ estimates noisy predictors of future $\beta^{-}$. (Average 1-year autocorrelation of $\beta^{-}$ is 0.435, but only 0.258 for highest volatility quintile stocks - Table 9A).
%         \item \textbf{Volatility Anomaly:} Stocks with very high total/idiosyncratic volatility tend to have anomalously low returns (Ang et al., 2006). This works against the high $\beta^{-} \rightarrow$ high return effect.
%     \end{enumerate}
%     \textbf{Excluding High Volatility Stocks (Table 9, Panel B):}
%     \begin{itemize}
%         \item Exclude the highest quintile of stocks based on past volatility (approx. 4\% of market cap).
%         \item Now, sorting the remaining stocks on past $\beta^{-}$ yields a \textbf{statistically significant positive relationship with future returns.} (High-Low spread of 0.34\% per month).
%     \end{itemize}
%     \textbf{Controlling for Volatility (Table 11 - Double Sorts):}
%     \begin{itemize}
%         \item Within each past volatility quintile (except the very highest), stocks with higher past $\beta^{-}$ have higher future size-and-book-to-market adjusted returns.
%         \item Averaging across volatility quintiles, the spread is 0.31\% per month.
%     \end{itemize}
% \end{frame}

\subsection{Bollerslev, Patton, and Quaedvlieg (2022)}

% Realized Semibetas Section
\begin{frame}{Defining Realized Semibetas}
  The traditional market beta ($\beta$) is decomposed as:
  $$
 \beta_t = \frac{\text{Cov}(r_t, f_t)}{\text{Var}(f_t)} 
$$
  This paper decomposes it into four components based on joint sign of asset return ($r_t$) and market return ($f_t$):
  \begin{itemize}
    \item $\beta^N$: Covariation when $r_t < 0$ and $f_t < 0$ (Both Negative)
    \item $\beta^P$: Covariation when $r_t > 0$ and $f_t > 0$ (Both Positive)
    \item $\beta^{M+}$: Covariation when $r_t < 0$ and $f_t > 0$ (Mixed: Market Positive, Asset Negative)
    \item $\beta^{M-}$: Covariation when $r_t > 0$ and $f_t < 0$ (Mixed: Market Negative, Asset Positive)
  \end{itemize}

  The paper defines mixed-sign semibetas as $\beta^{M+} = -M^+/\text{Var}(f)$ and $\beta^{M-} = -M^-/\text{Var}(f)$ for easier interpretation of risk premia.
  The decomposition is: $\beta = \beta^N + \beta^P - \beta^{M+} - \beta^{M-}$.
\end{frame}

\begin{frame}{Intuition}
  Imagine four assets, all with CAPM beta = 1:
  \begin{itemize}
    \item \textbf{Asset A (Normal):} Standard joint Normal distribution. Expected return benchmark.
    \item \textbf{Asset B ($\beta^N < \beta^P$):} Less correlated during market downturns. Desirable hedging property.
    \begin{itemize}
        \item \textbf{Predicted: Lower expected return} than Asset A.
    \end{itemize}
    \item \textbf{Asset C ($\beta^N > \beta^P$):} More correlated during market downturns. Undesirable from a mean-semivariance perspective.
    \begin{itemize}
        \item \textbf{Predicted: Higher expected return} than Asset A.
    \end{itemize}
    \item \textbf{Asset D ($\beta^N > \beta^P$, but $\beta^{M-} > \beta^{M+}$):} Also more correlated in downturns, but its mixed semicovariations offer superior hedging benefits (e.g., goes up more when market is down) compared to C.
    \begin{itemize}
        \item \textbf{Predicted: Lower expected return} than Asset C.
    \end{itemize}
  \end{itemize}
\end{frame}

% Empirical Findings
\begin{frame}{Pricing of Monthly Semibetas}
  Key findings from Fama-MacBeth regressions using monthly semibetas:
  \begin{itemize}
    \item \textbf{$\beta^N$ (both negative): Associated with significantly higher subsequent returns.}
    \begin{itemize}
        \item Investors demand a premium for assets that fall when the market falls.
    \end{itemize}
    \item \textbf{$\beta^{M-}$ (market negative, asset positive): Associated with significantly lower subsequent returns.}
    \begin{itemize}
        \item Assets that hedge market downturns (go up when the market is down) offer lower returns. This $\beta^{M-}$ is defined with a negative sign, so a positive coefficient means lower returns. The paper states, ``stocks with higher $\beta^{M-}$ are associated with significantly lower subsequent daily returns".
    \end{itemize}
    \item $\beta^P$ (both positive) and $\beta^{M+}$ (market positive, asset negative) appear \textbf{unpriced}.
    \item Results are robust to the inclusion of standard controls (size, B/M, momentum, IVOL, etc.).
    \item Semibeta model shows higher cross-sectional $R^2$ than traditional CAPM.
  \end{itemize}
\end{frame}

\section{Betting against beta}

\subsection{Frazzini and Pedersen (2014)}

\begin{frame}{Motivation: Leverage Constraints}
  \begin{itemize}
    \item Many investors, such as individuals, pension funds, and mutual funds, are \textbf{constrained in their use of leverage}.
    \item To achieve higher target returns, these constrained investors may \textbf{overweight high-beta assets} instead of leveraging low-beta assets.
    \item This behavior can lead to:
        \begin{itemize}
            \item High-beta assets becoming ``overpriced" (demanded by constrained investors), thus offering low alpha.
            \item Low-beta assets being relatively ``underpriced," offering higher alpha if leverage could be applied.
        \end{itemize}
    \item The paper builds on and extends the work of Black (1972) on CAPM with restricted borrowing.
  \end{itemize}
\end{frame}

% Theoretical Model Section
\begin{frame}{The Model}
  \begin{itemize}
    \item The paper presents a model with \textbf{heterogeneous investors} facing varying leverage and margin constraints.
    \item Some investors cannot use leverage, while others can but face margin requirements.
    \item These constraints ($m_i$) affect portfolio choices:
        \begin{itemize}
            \item Constrained investors who seek higher returns tend to tilt their portfolios towards high-beta assets.
            \item Unconstrained investors can exploit this by holding leveraged low-beta assets and shorting high-beta assets.
        \end{itemize}
    \item The model generates an equilibrium where funding constraints influence asset prices and lead to a \textbf{flatter SML}.
    \item Key variable: $\psi_t$, the average Lagrange multiplier, measuring the tightness of funding constraints.
  \end{itemize}
\end{frame}

\begin{frame}{Key Propositions from the Model}
  The model yields several empirically testable predictions:
  \begin{enumerate}
    \item \textbf{High beta is low alpha}:
        \begin{itemize}
            \item Alpha is inversely related to beta: $\alpha_s = \psi_t (1-\beta_s)$.
            \item The SML is flatter; intercept is $r_f + \psi_t$.
        \end{itemize}
    \item \textbf{Positive expected return of BAB factor}:
        \begin{itemize}
            \item A ``Betting Against Beta" (BAB) factor, which is long leveraged low-beta assets and short de-leveraged high-beta assets (market-neutral), produces significant positive risk-adjusted returns.
        \end{itemize}
    \item \textbf{Funding shocks and BAB returns}:
        \begin{itemize}
            \item When funding constraints tighten (e.g., margin requirements increase), the contemporaneous return of the BAB factor is low, and its future required return increases.
        \end{itemize}
    \item \textbf{Beta compression}:
        \begin{itemize}
            \item Increased funding liquidity risk compresses betas of all securities toward one.
        \end{itemize}
    \item \textbf{Constrained investors hold riskier assets}:
        \begin{itemize}
            \item More constrained investors hold higher-beta assets in their portfolios.
        \end{itemize}
  \end{enumerate}
\end{frame}

\begin{frame}{Methods}
    \begin{equation}
\hat{\beta}_i^{TS}=\hat{\rho} \frac{\hat{\sigma}_i}{\hat{\sigma}_m}
\end{equation}
\begin{itemize}
    \item a \textbf{one-year rolling standard deviation} for volatilities 
    \item a \textbf{five-year horizon for the correlation} to account for the fact that correlations appear to move more slowly than volatilities.
    \item one-day log returns to estimate volatilities
    \item overlapping three-day log returns, $r_{i, t}^{3 d}=\sum_{k=0}^2 \ln \left(1+r_{t+k}^i\right)$, for correlation to control for nonsynchronous trading (which affects only correlations).
\end{itemize}

\begin{equation}
\hat{\beta}_i=w_i \hat{\beta}_i^{T S}+\left(1-w_i\right) \hat{\beta}^{X S}
\end{equation}
Then shrink the time series estimate of beta $\hat{\beta}_i^{t s}$ toward the cross-sectional mean
\end{frame}

\begin{frame}{Following Literature}
\begin{itemize}
    \item Cederburg, Scott, and Michael S. O’Doherty, 2016, \textbf{Does it pay to bet against beta? On the conditional performance of the beta anomaly}, Journal of Finance 71, 737–774.
    \item Asness, Cliff, Andrea Frazzini, Niels Joachim Gormsen, and Lasse Heje Pedersen, 2020, \textbf{Betting against correlation: Testing theories of the low-risk effect}, Journal of Financial Economics 135, 629–652.
   \item Novy-Marx, Robert, and Mihail Velikov, 2022, \textbf{Betting against betting against beta}, Journal of Financial Economics 143, 80–106.
\end{itemize}
\end{frame}

\section{Focus on Different Times or Different Types Stocks}

\begin{frame}{Overview of Selected Papers}
  \begin{itemize}
    \item \textbf{Savor \& Wilson (2014):} Macroeconomic Announcement Days
    \item \textbf{Hong \& Sraer (2016):} Speculative Betas and Disagreement
    \item \textbf{Wang et al. (2017):} Reference-Dependent Preferences
    \item \textbf{Hendershott et al. (2020):} Night vs. Day Returns
    \item \textbf{Bodilsen et al. (2021):} FOMC Press Conferences
    \item \textbf{Chan \& Marsh (2022):} Leading Earnings Announcement Days
    \item \textbf{Huang et al. (2024):} Booms and Busts of Beta Arbitrage
  \end{itemize}
  \bigskip
  \textbf{Common Thread:} The relationship between beta and returns is not static but varies systematically with market conditions, information flow, and investor behavior.
\end{frame}

\begin{frame}{Savor \& Wilson (2014): A Tale of Two Days}

  \begin{itemize}
    \item \textbf{Focus:} Examines asset price behavior on days with scheduled macroeconomic news announcements (e.g., inflation, unemployment, FOMC) versus other days.
    \item \textbf{Key Findings:}
    \begin{itemize}
        \item Risky assets earn significantly higher average returns on announcement days (a-days).
        \item On a-days, stock market beta is \textbf{strongly positively related} to average returns.
        \item SML intercept is low/insignificant, slope close to market excess return on a-days.
        \item On non-announcement days (n-days), beta is unrelated or negatively related to returns.
    \end{itemize}
    \item \textbf{Implication for CAPM:} CAPM appears to hold conditionally during periods of important, pre-scheduled macroeconomic information release.
  \end{itemize}
\end{frame}

\begin{frame}{Hong \& Sraer (2016): Speculative Betas}

  \begin{itemize}
    \item \textbf{Focus:} Impact of investor disagreement about market prospects, especially with short-sales constraints.
    \item \textbf{Key Findings:}
    \begin{itemize}
        \item High-beta assets are more sensitive to aggregate disagreement, leading to greater divergence of opinion about their payoffs.
        \item With short-sales constraints, optimists drive up prices of high-beta assets.
        \item \textbf{High disagreement:} SML can be flat or even inverted (downward-sloping) due to overpricing of high-beta stocks.
        \item \textbf{Low disagreement:} SML tends to be upward-sloping (risk-sharing dominates).
    \end{itemize}
    \item \textbf{Implication for CAPM:} Behavioral factors (disagreement) and market frictions (short-sale constraints) can significantly distort the SML.
  \end{itemize}
\end{frame}

% --- Wang et al. (2017) ---
\begin{frame}{Wang et al. (2017): Reference-Dependent Preferences}
  \begin{itemize}
    \item \textbf{Focus:} How prior gains or losses (Capital Gains Overhang - CGO) affect the risk-return relationship, motivated by prospect theory.
    \item \textbf{Key Findings:}
    \begin{itemize}
        \item Among stocks with \textbf{high CGO} (investors in a gain position), the risk-return relation is \textbf{positive}.
        \item Among stocks with \textbf{low CGO} (investors in a loss position), the risk-return relation is \textbf{negative}.
        \item This holds for various risk measures (beta, volatility, idiosyncratic volatility).
    \end{itemize}
    \item \textbf{Implication for CAPM:} Investor psychology, specifically reference-dependent risk attitudes (risk-averse in gains, risk-seeking in losses), influences the observed SML.
  \end{itemize}
\end{frame}

\begin{frame}{Hendershott et al. (2020): A Tale of Night and Day}
  \begin{itemize}
    \item \textbf{Focus:} Differentiates between returns earned overnight (market closed) and during the trading day (market open).
    \item \textbf{Key Findings:}
    \begin{itemize}
        \item \textbf{Overnight:} Stock returns are \textbf{positively related} to beta (CAPM-like).
        \item \textbf{Trading Day:} Stock returns are \textbf{negatively related} to beta.
        \item The implied risk-free rate also differs significantly between night and day.
        \item Holds for US and international stocks, various portfolios.
    \end{itemize}
    \item \textbf{Implication for CAPM:} The failure of the 24-hour CAPM might be due to opposing SML slopes and different pricing regimes during trading vs. non-trading hours.
  \end{itemize}
\end{frame}

\begin{frame}{Bodilsen et al. (2021): FOMC Press Conferences}
  \begin{itemize}
    \item \textbf{Focus:} Differentiates FOMC announcement days \textit{with} a press conference (pc-days) from those \textit{without} (non-pc days) during 2011-2018.
    \item \textbf{Key Findings:}
    \begin{itemize}
        \item On \textbf{pc-days}, excess stock returns are strongly and \textbf{positively related} to beta.
        \item SML slope is positive and significant on pc-days.
        \item On non-pc days (and non-announcement days), the SML is flat or insignificant.
        \item Cross-sectional dispersion in betas declines substantially on pc-days.
    \end{itemize}
    \item \textbf{Implication for CAPM:} The \textit{way} information is communicated (not just the announcement itself) matters for asset pricing. PCs enhance the CAPM-like relationship.
  \end{itemize}
\end{frame}

\begin{frame}{Chan \& Marsh (2022): Earnings Announcement Days}
  \begin{itemize}
    \item \textbf{Focus:} Examines asset pricing on days when a cluster of influential S\&P 500 firms disclose quarterly earnings early in the season (LEADs).
    \item \textbf{Key Findings:}
    \begin{itemize}
        \item On \textbf{LEADs}, market betas have a strong and \textbf{positive relation} with average stock returns.
        \item SML intercept is insignificant on LEADs.
        \item On days other than LEADs, the beta-return relation is flat or negative.
        \item Market premium is significantly higher on LEADs.
    \end{itemize}
    \item \textbf{Implication for CAPM:} Concentrated, market-moving firm-specific news (early earnings from leaders) can create conditions where CAPM holds.
  \end{itemize}
\end{frame}

\begin{frame}{Huang et al. (2024): Booms and Busts of Beta Arbitrage}
  \begin{itemize}
    \item \textbf{Focus:} How arbitrage activity (measured by excess return comovement of beta-arbitrage stocks, CoBAR) affects the profitability and dynamics of beta-arbitrage strategies.
    \item \textbf{Key Findings:}
    \begin{itemize}
        \item \textbf{High CoBAR (high arbitrage activity):}
        \begin{itemize}
            \item Short-term: SML slopes downward (low-beta strategy profitable).
            \item Long-term (e.g., year 3): SML slopes upward (reversal, strategy loses money) $\implies$ overshooting.
        \end{itemize}
        \item \textbf{Low CoBAR (low arbitrage activity):}
        \begin{itemize}
            \item Short-term: SML weakly upward sloping.
            \item Long-term: SML slopes downward (delayed correction).
        \end{itemize}
        \item Proposes a positive-feedback channel via firm leverage.
    \end{itemize}
    \item \textbf{Implication for CAPM:} Arbitrage activity itself creates dynamic patterns in the SML, leading to periods of apparent CAPM failure and correction/overcorrection.
  \end{itemize}
\end{frame}

\begin{frame}{Common Themes and Insights}
  \begin{itemize}
    \item \textbf{Conditional Validity of CAPM:} The unconditional CAPM generally fails, but beta appears to be priced positively under specific, identifiable conditions.
    \item \textbf{Role of Information Flow:}
    \begin{itemize}
        \item CAPM performs better during periods of significant, market-wide information release (macro news, FOMC PCs, LEADs).
        \item The \textit{nature} and \textit{intensity} of information (and its communication) matter.
    \end{itemize}
    \item \textbf{Investor Behavior and Market Frictions:}
    \begin{itemize}
        \item Disagreement (Hong \& Sraer) and reference-dependent preferences (Wang et al.) can significantly alter the SML.
        \item Arbitrage activity and its limits (Huang et al.) lead to dynamic patterns, not always towards efficiency.
    \end{itemize}
    \item \textbf{Market Structure and Timing:}
    \begin{itemize}
        \item Pricing dynamics differ drastically between trading hours and overnight (Hendershott et al.).
    \end{itemize}
    \item \textbf{Beta is Not (Completely) Dead:} While not a universal explainer of returns, beta remains a relevant measure of risk, especially when conditioned on the right state variables.
  \end{itemize}
\end{frame}

\section{Puzzles and Future Directions}

\begin{frame}{Remaining Puzzles and Future Research}
  \begin{itemize}
    \item \textbf{Unifying Theory:} Is there an overarching framework that can explain \textit{why} these specific conditions (announcements, sentiment, CGO, time of day) matter for beta's pricing?
    \item \textbf{Interaction Effects:} How do these different conditionalities interact? For example, do macro announcements during high disagreement periods show a different SML than during low disagreement?
    \item \textbf{Mechanism Identification:}
    \begin{itemize}
        \item What are the precise micro-foundations? Changes in risk aversion, perceived risk, liquidity, attention, or arbitrage capital?
        \item Why does beta compression occur on some announcement days (e.g., Bodilsen et al.)?
    \end{itemize}
    \item \textbf{Predictability and Arbitrage:}
    \begin{itemize}
        \item Can these conditional CAPM findings lead to robust, out-of-sample profitable trading strategies after accounting for transaction costs and implementation challenges?
        \item How do sophisticated arbitrageurs respond to these predictable patterns (as partly addressed by Huang et al.)?
    \end{itemize}
    \item \textbf{International Evidence:} While some papers provide international evidence (e.g., Hendershott et al.), more is needed for other conditionalities.
  \end{itemize}
\end{frame}


\end{document}